\documentclass[11pt,a4paper]{article}

\usepackage[T1]{fontenc}
\usepackage[utf8]{inputenc}
\usepackage{geometry}
\usepackage{array}
\usepackage{longtable}
\usepackage{amsmath}
\usepackage{amssymb}
\usepackage{url}

\geometry{margin=22mm}
\newcommand{\code}[1]{\texttt{#1}}
\newcommand{\unit}[1]{\,\mathrm{#1}}
\newcommand{\NA}{--}
\newcommand{\toprule}{\hline}
\newcommand{\midrule}{\hline}
\newcommand{\bottomrule}{\hline}
\newcolumntype{L}[1]{>{\raggedright\arraybackslash}p{#1}}

\title{CitySim Environmental Models\\
       Pollutant Dispersion and Outdoor Noise Propagation}
\author{Implementation recap and XML input reference}
\date{\today}

\begin{document}
\maketitle

\begin{abstract}
This note documents the pollutant-dispersion and outdoor-noise modules added to
CitySim. Both modules are evaluated hourly, accept multiple point emitters and
multiple point receptors, use the CitySim three-dimensional coordinate system,
and write dedicated tab-separated result files. The pollutant module is a
Gaussian-plume screening model inspired by the use case covered by Soft-P's
AERO software. It is not an implementation of the complete Polish regulatory
AERO methodology. The noise module is a single-frequency engineering subset of
ISO~9613-1 and ISO~9613-2:2024. It is not a certified full implementation of
ISO~9613-2:2024. The scope and limitations are stated explicitly below.
\end{abstract}

\tableofcontents

\section{Integration in CitySim}

The two new model tags are children of the CitySim \code{<District>} element:
\code{<Pollutant>} and \code{<Noise>}. More than one model tag can be declared.
Each tag contains one or several \code{<PointEmittor>} elements and one or
several \code{<PointReceptor>} elements. The spelling \code{Emittor}, with two
letters ``t'', is the XML spelling currently implemented in the source code.

The models are computed after the regular CitySim hourly calculations. They use
the same climate object and the same coordinate frame as the rest of the scene.
The result files are:
\begin{itemize}
  \item \code{<input-base-name>\_Pollutant.tsv};
  \item \code{<input-base-name>\_Noise.tsv}.
\end{itemize}

Both files use the usual CitySim convention: one header line beginning with
\code{\#timeStep}, followed by one line per simulated hour. Identifiers and
optional keys appear in column names so that several models, receptors, and
emitters can be distinguished.

\section{Pollutant Dispersion Model}

\subsection{Physical scope}

The implemented pollutant model is a steady-state Gaussian plume evaluated at
each CitySim hourly time step. It is suitable for screening calculations from
point sources such as chimneys or discretised road-traffic emissions. The
hourly climate data determine the wind direction, the default reference wind
speed, and the automatically selected atmospheric stability class.

The model uses the CitySim coordinates of sources and receptors. Building
geometry does \emph{not} modify pollutant dispersion: there is no street-canyon
flow, facade interception, recirculation zone, or computational-fluid-dynamics
calculation. The geometry is presently a shared positioning frame only.

Soft-P's AERO application calculates atmospheric pollutant dispersion according
to the Polish regulatory framework \cite{softp-aero,polish-regulation}. The
present implementation covers a simpler Gaussian-plume use case and should not
be presented as a reproduction of the full AERO methodology.

\subsection{Wind coordinates and logarithmic wind profile}

Meteorological wind direction $\theta$ is the direction \emph{from which} the
wind blows. It is converted to a downwind unit vector:
\begin{equation}
  w_x = -\sin(\theta), \qquad w_y = -\cos(\theta).
\end{equation}
For an emitter--receptor horizontal displacement
$(\Delta x,\Delta y)=(x_r-x_e,y_r-y_e)$, the downwind and crosswind coordinates
are:
\begin{align}
  x &= \Delta x\,w_x + \Delta y\,w_y, \\
  y &= -\Delta x\,w_y + \Delta y\,w_x.
\end{align}
Upwind receptors ($x\leq0$) receive no contribution from that emitter. For
numerical stability, positive downwind distances are clamped to
$x\geq x_{\min}$, where $x_{\min}$ is configured by \code{downwindMin}.

The effective emitter height is:
\begin{equation}
  H = z_e + \Delta h_{\mathrm{plume}},
\end{equation}
where $z_e$ is the emitter coordinate and $\Delta h_{\mathrm{plume}}$ is the
user-entered \code{plumeRise}. The wind speed at this effective height follows
a logarithmic profile:
\begin{equation}
  u(H) =
  \max\left[
    u_{\min},
    u_{\mathrm{ref}}
    \frac{\ln(H^\ast/z_0)}{\ln(z_{\mathrm{ref}}/z_0)}
  \right],
  \qquad
  H^\ast=\max(H,1.01z_0),
  \label{eq:pollutant-wind-profile}
\end{equation}
where $z_0$ is the roughness length, $z_{\mathrm{ref}}$ is the reference
height, and $u_{\min}$ avoids singular low-wind calculations. By default,
$u_{\mathrm{ref}}$ is the hourly CitySim climate wind speed. The optional
\code{windSpeed} XML attribute overrides it.

This logarithmic profile is a CitySim model choice. It is not the
stability-dependent power-law wind profile used by the complete AERO regulatory
workflow.

\subsection{Gaussian plume concentration}

For one emitter $e$ and one receptor $r$, the concentration contribution is:
\begin{align}
 C_{e,r} &=
 \frac{Q_e}{2\pi u(H)\sigma_y\sigma_z}
 \exp\left[-\frac{1}{2}\left(\frac{y}{\sigma_y}\right)^2\right]
 V_z F_{\mathrm{decay}}F_{\mathrm{dep}}, \\
 V_z &=
 \exp\left[-\frac{1}{2}\left(\frac{z_r-H}{\sigma_z}\right)^2\right]
 + I_g
 \exp\left[-\frac{1}{2}\left(\frac{z_r+H}{\sigma_z}\right)^2\right],
 \label{eq:gaussian-plume}
\end{align}
where:
\begin{itemize}
  \item $Q_e$ is the emitter mass flow rate in $\unit{kg\,s^{-1}}$;
  \item $u(H)$ is the wind speed at the effective emitter height;
  \item $\sigma_y$ and $\sigma_z$ are horizontal and vertical plume spreads;
  \item $z_r$ is the receptor height;
  \item $I_g=1$ when \code{groundReflection=true}, otherwise $I_g=0$.
\end{itemize}

The optional first-order attenuation factors are:
\begin{align}
 F_{\mathrm{decay}} &= \exp\left(-k\frac{x}{u(H)}\right), \\
 F_{\mathrm{dep}} &=
 \exp\left[
   -\frac{v_d x}{u(H)\max(h_{\mathrm{mix}},1\unit{m})}
 \right],
 \label{eq:pollutant-optional-losses}
\end{align}
where $k$ is \code{decayRate}, $v_d$ is \code{depositionVelocity}, and
$h_{\mathrm{mix}}$ is \code{mixingHeight}. These exponential factors are
screening approximations. They are not the complete AERO suspended-dust and
deposition calculations.

Contributions from all emitters are summed linearly with an optional background
concentration:
\begin{equation}
  C_r = C_{\mathrm{background}} + \sum_e C_{e,r}.
\end{equation}

\subsection{Pasquill--Gifford plume spreads}

Unless both \code{sigmaY} and \code{sigmaZ} are given manually, the module uses
the following Pasquill--Gifford screening coefficients \cite{turner-workbook}.
Distances and plume spreads are expressed in metres.

\begin{center}
\small
\begin{tabular}{c l l}
\toprule
Class & $\sigma_y(x)$ & $\sigma_z(x)$ \\
\midrule
A & $0.22x(1+0.0001x)^{-1/2}$ & $0.20x$ \\
B & $0.16x(1+0.0001x)^{-1/2}$ & $0.12x$ \\
C & $0.11x(1+0.0001x)^{-1/2}$ & $0.08x(1+0.0002x)^{-1/2}$ \\
D & $0.08x(1+0.0001x)^{-1/2}$ & $0.06x(1+0.0015x)^{-1/2}$ \\
E & $0.06x(1+0.0001x)^{-1/2}$ & $0.03x(1+0.0003x)^{-1}$ \\
F & $0.04x(1+0.0001x)^{-1/2}$ & $0.016x(1+0.0003x)^{-1}$ \\
\bottomrule
\end{tabular}
\end{center}

The horizontal spread is multiplied by \code{sigmaYFactor}. The vertical
spread is multiplied by \code{sigmaZFactor}. When \code{stabilityClass} is
\code{auto}, CitySim selects a conventional class from hourly global horizontal
irradiance, cloud-cover fraction, and wind speed. A manual class \code{A} to
\code{F} can be provided for controlled studies.

\begin{center}
\small
\begin{tabular}{l l l}
\toprule
Condition & Wind condition & Selected class \\
\midrule
$I_{gh}>600\unit{W\,m^{-2}}$ & $u<2$, $2\leq u<3$, $3\leq u<5$, $u\geq5$ & A, B, C, D \\
$300<I_{gh}\leq600\unit{W\,m^{-2}}$ & $u<2$, $2\leq u<5$, $u\geq5$ & B, C, D \\
$10<I_{gh}\leq300\unit{W\,m^{-2}}$ & $u<2$, $u\geq2$ & C, D \\
$I_{gh}\leq10\unit{W\,m^{-2}}$, cloud cover $>0.5$ & any & D \\
$I_{gh}\leq10\unit{W\,m^{-2}}$, cloud cover $\leq0.5$ & $u<2$, $2\leq u<5$, $u\geq5$ & F, E, D \\
\bottomrule
\end{tabular}
\end{center}

\subsection{Pollutant XML parameters}

\begin{longtable}{>{\ttfamily\raggedright\arraybackslash}p{0.29\linewidth} L{0.12\linewidth} L{0.11\linewidth} L{0.34\linewidth}}
\toprule
\normalfont Attribute & Unit & Default & Meaning \\
\midrule
\endhead
\multicolumn{4}{l}{\textbf{\code{<Pollutant>} attributes}} \\
model & -- & plume & Model identifier. The implemented calculation is the plume model. \\
id & -- & Pollutant & Result-column model identifier. \\
key & -- & model & Optional result-column key, often used for the species name such as NOx or PM10. \\
windSpeed & $\unit{m\,s^{-1}}$ & climate & Optional manual wind-speed override at \code{windReferenceHeight}. \\
windReferenceHeight & $\unit{m}$ & 10 & Reference height $z_{\mathrm{ref}}$ for the logarithmic profile. \\
roughnessLength & $\unit{m}$ & 0.1 & Aerodynamic roughness length $z_0$; must be positive. \\
windSpeedMin & $\unit{m\,s^{-1}}$ & 0.5 & Minimum wind speed used by the plume equation; must be positive. \\
downwindMin & $\unit{m}$ & 1 & Minimum positive downwind distance used numerically. \\
backgroundConcentration & $\unit{kg\,m^{-3}}$ & 0 & Concentration added at every receptor. \\
groundReflection & boolean & true & Include the reflected Gaussian image-source term. \\
stabilityClass & -- & auto & Automatic selection or manual class A--F. \\
sigmaY & $\unit{m}$ & \NA & Manual constant $\sigma_y$; active only when \code{sigmaZ} is also set. \\
sigmaZ & $\unit{m}$ & \NA & Manual constant $\sigma_z$; active only when \code{sigmaY} is also set. \\
sigmaYFactor & -- & 1 & Multiplier applied to $\sigma_y$. \\
sigmaZFactor & -- & 1 & Multiplier applied to $\sigma_z$. \\
decayRate & $\unit{s^{-1}}$ & 0 & Optional first-order pollutant decay rate $k$. \\
depositionVelocity & $\unit{m\,s^{-1}}$ & 0 & Optional simplified deposition velocity $v_d$. \\
mixingHeight & $\unit{m}$ & 1000 & Mixing height used by the simplified deposition term. \\
\midrule
\multicolumn{4}{l}{\textbf{\code{<PointEmittor>} attributes inside \code{<Pollutant>}}} \\
id & -- & sequence & Emitter identifier. \\
key & -- & empty & Optional emitter key. \\
x, y, z & $\unit{m}$ & 0 & Emitter position in CitySim scene coordinates. \\
emissionRate & $\unit{kg\,s^{-1}}$ & 0 & Pollutant mass flow rate $Q_e$. \\
plumeRise & $\unit{m}$ & 0 & User-entered effective plume rise above the source point. \\
\midrule
\multicolumn{4}{l}{\textbf{\code{<PointReceptor>} attributes inside \code{<Pollutant>}}} \\
id & -- & sequence & Receptor identifier. \\
key & -- & empty & Optional receptor key. \\
x, y, z & $\unit{m}$ & 0 & Receptor position in CitySim scene coordinates. \\
\bottomrule
\end{longtable}

\subsection{Pollutant XML example}

Declare one \code{<Pollutant>} block per species when simulating several
pollutants such as NOx, PM2.5, and PM10. Each species can have different
emission rates, background concentrations, decay rates, and deposition
velocities.

\begin{verbatim}
<District ...>
  <Pollutant id="NOx" key="NOx" model="plume"
             windReferenceHeight="10"
             roughnessLength="0.1"
             windSpeedMin="0.5"
             backgroundConcentration="0">
    <PointEmittor id="chimney_1" key="chimney"
                  x="130" y="240" z="32"
                  plumeRise="5"
                  emissionRate="0.001"/>
    <PointEmittor id="traffic_1" key="traffic"
                  x="95" y="180" z="0.5"
                  emissionRate="0.00002"/>
    <PointReceptor id="street_1" x="150" y="170" z="1.5"/>
    <PointReceptor id="roof_1" x="210" y="120" z="28"/>
  </Pollutant>
</District>
\end{verbatim}

To force a fixed reference wind speed for a sensitivity calculation, set
\code{windSpeed} to \code{4}, for example. When omitted, the CitySim climate wind speed is
used hourly.

\subsection{Pollutant result columns}

The pollutant TSV file contains:
\begin{itemize}
  \item mean effective emitter-height \code{windSpeed(m/s)} for the model;
  \item hourly climate \code{windDirection(deg)};
  \item selected \code{stabilityClass};
  \item total \code{concentration(kg/m3)} at each receptor;
  \item each emitter's \code{concentration(kg/m3)} contribution at each receptor.
\end{itemize}
The concentration equation itself uses the individual wind speed evaluated at
each emitter height. The reported wind speed is their arithmetic mean.

\section{Outdoor Noise Propagation Model}

\subsection{Physical scope}

The noise module estimates steady-state outdoor sound-pressure levels from
point emitters to point receptors. Unlike the pollutant module, it actively
uses CitySim three-dimensional geometry. Building walls and roofs are collected
from the district, tested for ray obstruction, used to construct top-edge
diffraction paths, and used for approximate first-order wall reflections.

The implemented attenuation terms follow the engineering structure of
ISO~9613-2:2024 \cite{iso9613-2}, with atmospheric absorption derived from
ISO~9613-1 \cite{iso9613-1}. The implementation currently evaluates one
representative frequency rather than aggregating octave bands. It is an
ISO-inspired CitySim engineering subset, not a certified full
ISO~9613-2:2024 implementation.

\subsection{Energy aggregation}

For one propagation path $p$ between an emitter and a receptor:
\begin{equation}
  L_{p} = L_W + D_c - A_{p},
\end{equation}
where $L_W$ is the point-source sound power level in
$\unit{dB}$ re $1\unit{pW}$, $D_c$ is the directivity correction, and $A_p$ is
the sum of path-specific attenuation terms.

Direct, diffracted, and reflected paths are combined incoherently:
\begin{equation}
  L_{e,r} = 10\log_{10}\left(\sum_p 10^{L_p/10}\right).
\end{equation}
Emitter contributions and the optional background level are also combined as
energies:
\begin{equation}
  L_r = 10\log_{10}\left(
    10^{L_{\mathrm{background}}/10}
    + \sum_e 10^{L_{e,r}/10}
  \right).
\end{equation}

\subsection{Direct-path attenuation}

For an unobstructed source--receptor path:
\begin{equation}
  A_p = A_{\mathrm{div}} + A_{\mathrm{atm}} + A_{\mathrm{gr}}.
\end{equation}
Spherical spreading is:
\begin{equation}
  A_{\mathrm{div}} = 20\log_{10}\left[\max(d,0.1\unit{m})\right]+11,
  \label{eq:noise-div}
\end{equation}
where $d$ is the three-dimensional source--receptor distance.

Atmospheric absorption is:
\begin{equation}
  A_{\mathrm{atm}} = \alpha\frac{d}{1000},
  \label{eq:noise-atm}
\end{equation}
where $\alpha$ is expressed in $\unit{dB\,km^{-1}}$. If it is not manually
specified, $\alpha$ is computed hourly from CitySim temperature, relative
humidity, and atmospheric pressure according to ISO~9613-1.

The speed of sound and wavelength used for diffraction are:
\begin{equation}
  c = 331.3\sqrt{1+\frac{T_C}{273.15}},
  \qquad
  \lambda=\frac{c}{f}.
\end{equation}

\subsection{ISO 9613-1 atmospheric absorption}

The ISO~9613-1 coefficient is implemented explicitly. Let $T$ be the absolute
temperature, $T_0=293.15\unit{K}$, $T_{\mathrm{tp}}=273.16\unit{K}$,
$p_r=p/101325$, and $h$ the molar humidity:
\begin{align}
  \frac{p_{\mathrm{sat}}}{p_0}
    &= 10^{-6.8346(T_{\mathrm{tp}}/T)^{1.261}+4.6151}, \\
  h &= 100\,RH\,\frac{p_{\mathrm{sat}}/p_0}{p_r}, \\
  f_{\mathrm{r,O}} &=
    \frac{24+4.04\times10^4 h(0.02+h)/(0.391+h)}{p_r}, \\
  f_{\mathrm{r,N}} &=
    \frac{(T/T_0)^{-1/2}
    \left[9+280h\exp\left(-4.17((T/T_0)^{-1/3}-1)\right)\right]}{p_r}.
\end{align}
For selected frequency $f$, the hourly absorption coefficient is:
\begin{align}
 \alpha ={}& 1000(8.686)f^2
 \left\{
 \frac{1.84\times10^{-11}}{p_r}\sqrt{\frac{T}{T_0}}
 + \left(\frac{T}{T_0}\right)^{-5/2}
 \left[
 \frac{0.01275e^{-2239.1/T}}
      {f_{\mathrm{r,O}}+f^2/f_{\mathrm{r,O}}}
 \right.\right. \nonumber\\
 &\left.\left.\hspace{38mm}
 + \frac{0.1068e^{-3352/T}}
      {f_{\mathrm{r,N}}+f^2/f_{\mathrm{r,N}}}
 \right]
 \right\}.
 \label{eq:noise-alpha}
\end{align}

\subsection{ISO 9613-2 ground effect}

The selected frequency is mapped to the closest nominal octave band:
$63$, $125$, $250$, $500$, $1000$, $2000$, $4000$, or $8000\unit{Hz}$.
The implementation uses one global ground factor $G$ for source, middle, and
receptor regions:
\begin{itemize}
  \item $G=0$: acoustically hard ground, such as paving or water;
  \item $G=1$: porous ground, such as grass or soil;
  \item $0<G<1$: mixed-ground screening approximation.
\end{itemize}

For horizontal source--receptor distance $d_p$ and acoustic heights $h_s$ and
$h_r$, define:
\begin{align}
  q &= \max\left(0,1-\frac{30(h_s+h_r)}{d_p}\right), \\
  K_{\mathrm{geo}} &=
  \frac{d_p^2+(h_s-h_r)^2}{d_p^2+(h_s+h_r)^2}.
\end{align}
The ground attenuation is:
\begin{equation}
  A_{\mathrm{gr}} =
  -10\log_{10}\left[
    1 + \left(10^{-A'_{\mathrm{gr}}/10}-1\right)K_{\mathrm{geo}}
  \right],
  \qquad
  A'_{\mathrm{gr}}=A_s+A_r+A_m.
  \label{eq:noise-ground}
\end{equation}
The source and receptor terms use the same function,
$A_s=A_G(h_s)$ and $A_r=A_G(h_r)$. The middle-region term is
$A_m=-3q$ at $63\unit{Hz}$ and $A_m=-3q(1-G)$ otherwise. The implemented
$A_G(h)$ functions are:

\begin{center}
\scriptsize
\begin{tabular}{c p{0.76\linewidth}}
\toprule
Band & $A_G(h)$ \\
\midrule
$63$ &
$-1.5$ \\
$125$ &
$-1.5 + G\{1.5+3e^{-0.12(h-5)^2}(1-e^{-d_p/50})
+5.7e^{-0.09h^2}(1-e^{-2.8\times10^{-6}d_p^2})\}$ \\
$250$ &
$-1.5 + G\{1.5+8.6e^{-0.09h^2}(1-e^{-d_p/50})\}$ \\
$500$ &
$-1.5 + G\{1.5+14e^{-0.46h^2}(1-e^{-d_p/50})\}$ \\
$1000$ &
$-1.5 + G\{1.5+5e^{-0.9h^2}(1-e^{-d_p/50})\}$ \\
$2000$, $4000$, $8000$ &
$-1.5(1-G)$ \\
\bottomrule
\end{tabular}
\end{center}

\subsection{Top-edge diffraction over buildings}

If the direct ray is blocked, CitySim constructs a screened path over the
intersected building top edges. Hidden low edges are pruned using an upper
rubber-band path. Each resulting segment is checked again against the scene.
If no clear top-edge path can be constructed, the direct/diffracted
contribution is omitted rather than allowing sound to pass through a building.

For a valid screened path, let:
\begin{itemize}
  \item $d$ be the direct source--receptor distance;
  \item $d_{\mathrm{dif}}$ be the sum of all screened-path segment lengths;
  \item $z=d_{\mathrm{dif}}-d$ be the path-length difference;
  \item $e$ be the distance along intermediate edges between the first and last
        diffraction edges;
  \item $d_{\mathrm{SS}}$ and $d_{\mathrm{SR}}$ be the first and last screened
        segment lengths.
\end{itemize}
The implemented ISO~9613-2:2024 top-edge attenuation uses:
\begin{align}
  C_2 &= 20, \\
  C_3 &=
  \begin{cases}
    1, & \text{one edge}, \\
    \dfrac{1+(5\lambda/e)^2}{1/3+(5\lambda/e)^2}, & \text{multiple edges},
  \end{cases} \\
  z_{\min} &= -\frac{2\lambda}{C_2C_3}, \\
  K_{\mathrm{met}} &=
  \exp\left[
    -\frac{1}{2000}
    \sqrt{
      \frac{
        [\max(d_{\mathrm{SS}},d_{\mathrm{SR}})+e]
        \min(d_{\mathrm{SS}},d_{\mathrm{SR}})d
      }{
        2(z-z_{\min})
      }
    }
  \right], \\
  A_{\mathrm{bar}} &=
  10\log_{10}\left[
    1+
    \left(2+\frac{C_2}{\lambda}C_3z\right)K_{\mathrm{met}}
  \right].
  \label{eq:noise-barrier}
\end{align}
When $z\leq z_{\min}$, $A_{\mathrm{bar}}=0$. For a screened ray, spreading is
evaluated at direct distance $d$, atmospheric absorption is evaluated along the
screened path length, and the applied ground term is
$\min(A_{\mathrm{gr}},0)$.

\subsection{Approximate first-order wall reflections}

When \code{reflectionOrder="1"}, vertical walls are considered for one
image-source reflection. Source and receptor must lie on the same side of the
reflecting wall. Both segments from source to reflection point and from
reflection point to receptor must be unobstructed.

For wall reflected-energy coefficient $\rho_w$:
\begin{equation}
  A_{\mathrm{refl}}=-10\log_{10}(\rho_w).
  \label{eq:noise-reflection}
\end{equation}
When only an absorption coefficient $\alpha_w$ is supplied in the wall
composition, CitySim uses $\rho_w=1-\alpha_w$.
A reflected path uses:
\begin{equation}
  A_p=A_{\mathrm{div}}(d_{\mathrm{refl}})
     +A_{\mathrm{atm}}(d_{\mathrm{refl}})
     +A_{\mathrm{refl}}.
\end{equation}

\subsection{Noise XML parameters}

\begin{longtable}{>{\ttfamily\raggedright\arraybackslash}p{0.29\linewidth} L{0.12\linewidth} L{0.11\linewidth} L{0.34\linewidth}}
\toprule
\normalfont Attribute & Unit & Default & Meaning \\
\midrule
\endhead
\multicolumn{4}{l}{\textbf{\code{<Noise>} attributes}} \\
id & -- & Noise & Result-column model identifier. \\
key & -- & id & Optional result-column key. \\
frequency & $\unit{Hz}$ & 500 & Representative frequency used for all attenuation terms. \\
atmosphericAlpha & $\unit{dB\,km^{-1}}$ & climate & Optional manual override. When omitted, use hourly ISO~9613-1 absorption from CitySim climate data. \\
groundFactor & -- & 0 & Global $G$: 0 for hard ground, 1 for porous ground. \\
wallAbsorptionCoeff & -- & 0.1 & Fallback absorbed-energy fraction for walls whose composition has no acoustic coefficient; range 0--1. \\
wallReflectivityCoeff & -- & 0.9 & Fallback reflected-energy fraction. If omitted, it is one minus \code{wallAbsorptionCoeff}. \\
reflectionOrder & -- & 1 & 0 for direct/diffracted propagation only; 1 to include first-order wall reflections. \\
backgroundLevel & $\unit{dB}$ & -300 & Incoherent background level. Values at or below $-300$ disable it. \\
directivityCorrection & $\unit{dB}$ & 0 & Default source directivity correction $D_c$. \\
\midrule
\multicolumn{4}{l}{\textbf{Acoustic wall-composition attributes on \code{<WallType>} or \code{<Composite>}}} \\
AcousticAbsorptionCoeff & -- & 0.1 & Absorbed-energy fraction for walls using this composition. If reflectivity is omitted, reflectivity becomes $1-\alpha_w$. \\
AcousticReflectivityCoeff & -- & 0.9 & Reflected-energy fraction for walls using this composition. This coefficient is used by first-order reflections. \\
\multicolumn{4}{l}{\code{AbsorptionCoeff} and \code{ReflectivityCoeff} are accepted aliases.} \\
\midrule
\multicolumn{4}{l}{\textbf{\code{<PointEmittor>} attributes inside \code{<Noise>}}} \\
id & -- & sequence & Emitter identifier. \\
key & -- & empty & Optional emitter key. \\
x, y, z & $\unit{m}$ & 0 & Emitter position in CitySim scene coordinates. \\
heightAboveGround & $\unit{m}$ & z & Acoustic source height above local ground used by $A_{\mathrm{gr}}$. \\
powerLevel & $\unit{dB}$ & 0 & Sound power level $L_W$, re $1\unit{pW}$. \\
directivityCorrection & $\unit{dB}$ & model value & Optional emitter-specific override of $D_c$. \\
\midrule
\multicolumn{4}{l}{\textbf{\code{<PointReceptor>} attributes inside \code{<Noise>}}} \\
id & -- & sequence & Receptor identifier. \\
key & -- & empty & Optional receptor key. \\
x, y, z & $\unit{m}$ & 0 & Receptor position in CitySim scene coordinates. \\
heightAboveGround & $\unit{m}$ & z & Acoustic receptor height above local ground used by $A_{\mathrm{gr}}$. \\
\bottomrule
\end{longtable}

\subsection{Noise XML example}

\begin{verbatim}
<District ...>
  <WallType id="1" name="ConcreteFacade" category="wall"
            AbsorptionCoeff="0.1"
            ReflectivityCoeff="0.9">
    ...
  </WallType>
  <Noise id="RoadNoise" key="road" frequency="500"
         groundFactor="0.5"
         reflectionOrder="1"
         backgroundLevel="30">
    <PointEmittor id="road_1"
                  x="110" y="180" z="0.5"
                  heightAboveGround="0.5"
                  powerLevel="90"
                  directivityCorrection="0"/>
    <PointReceptor id="facade_1"
                   x="185" y="160" z="4"
                   heightAboveGround="4"/>
  </Noise>
</District>
\end{verbatim}

The \code{z} coordinate places a point in the CitySim scene. The optional
\code{heightAboveGround} attribute supplies the local acoustic height used by
the ground-effect model. Set it explicitly for elevated terrain and roof-level
points. Normally omit \code{atmosphericAlpha}; it remains available as a
manual override for calibrated or controlled studies.

\subsection{Noise result columns}

The noise TSV file contains:
\begin{itemize}
  \item hourly \code{atmosphericAlpha(dB/km)};
  \item total \code{level(dB)} at each receptor;
  \item each emitter's \code{level(dB)} contribution at each receptor.
\end{itemize}

\section{Model Comparison and Limitations}

\begin{longtable}{L{0.24\linewidth} L{0.34\linewidth} L{0.34\linewidth}}
\toprule
Aspect & Pollutant module & Noise module \\
\midrule
\endhead
Primary purpose &
Hourly Gaussian-plume concentration screening. &
Hourly outdoor sound-propagation screening. \\
CitySim climate use &
Wind speed, wind direction, irradiance, and cloud cover. &
Temperature, relative humidity, and atmospheric pressure. \\
CitySim geometry use &
Coordinates only. Buildings do not alter the plume. &
Coordinates, obstruction tests, top-edge screening, and first-order wall reflections. \\
Multiple emitters and receptors &
Yes. Mass concentrations are summed linearly. &
Yes. Sound levels are summed as incoherent energies. \\
Main reference family &
Gaussian-plume screening literature and AERO use case. &
ISO~9613-1 and ISO~9613-2:2024 engineering structure. \\
\bottomrule
\end{longtable}

\subsection{Pollutant limitations}

\begin{itemize}
  \item No building-wake, urban-canyon, or terrain-flow correction.
  \item No stack-physics plume-rise calculation from exit velocity, stack
        diameter, or exhaust temperature. \code{plumeRise} is entered manually.
  \item No complete Polish AERO regulatory meteorology, coefficients, or
        suspended-dust and deposition method.
  \item No chemical reaction network. \code{decayRate} is a single simplified
        first-order loss.
  \item No line-source primitive. Traffic roads are represented by several
        point emitters.
\end{itemize}

\subsection{Noise limitations}

\begin{itemize}
  \item One representative frequency is evaluated. Full octave-band source
        spectra and A-weighted band aggregation are not yet implemented.
  \item Diffraction paths go over building tops. Lateral paths around facade
        sides are not evaluated and compared.
  \item One global \code{groundFactor} is used instead of spatial ground
        regions mapped from the CitySim scene.
  \item Reflections are approximate first-order image-source paths. Higher
        orders and more complete object-reflection handling are not included.
  \item Vegetation, miscellaneous-object attenuation, line sources, and area
        sources are not included.
  \item Complex non-convex surfaces remain an engineering approximation because
        polygons are fan-triangulated for ray tests.
\end{itemize}

\section{Recommended Configuration Practice}

\begin{enumerate}
  \item Use one \code{<Pollutant>} block per chemical species. Set
        species-specific emissions for NOx, PM2.5, and PM10.
  \item Use CitySim climate wind speed by default. Apply \code{windSpeed} only
        for controlled sensitivity studies.
  \item Calibrate \code{plumeRise}, background concentration, roughness, and
        optional pollutant losses against the intended study scale.
  \item For noise, supply octave-specific runs when spectral sensitivity
        matters: use one \code{<Noise>} block per representative octave band
        until native octave aggregation is implemented.
  \item Set noise \code{heightAboveGround} explicitly whenever scene $z$ is not
        the local acoustic height.
  \item Treat both modules as transparent engineering-screening tools. Use a
        regulatory implementation when legal compliance must be demonstrated.
\end{enumerate}

\section{References}

\begin{thebibliography}{9}

\bibitem{softp-aero}
Soft-P,
\emph{AERO 2010 -- calculation of air pollutant dispersion},
software description.
\url{https://soft-p.com.pl/en/software/aero-2010/}

\bibitem{polish-regulation}
Minister of the Environment of the Republic of Poland,
\emph{Regulation of 26 January 2010 on reference values for certain substances
in the air},
Dz.U. 2010 nr 16 poz. 87.
\url{https://isap.sejm.gov.pl/isap.nsf/DocDetails.xsp?id=WDU20100160087}

\bibitem{turner-workbook}
D. B. Turner,
\emph{Workbook of Atmospheric Dispersion Estimates},
U.S. Environmental Protection Agency, AP-26, revised edition, 1970.
\url{https://nepis.epa.gov/Exe/ZyPURL.cgi?Dockey=2000JCID.TXT}

\bibitem{iso9613-1}
International Organization for Standardization,
\emph{ISO 9613-1:1993, Acoustics -- Attenuation of sound during propagation
outdoors -- Part 1: Calculation of the absorption of sound by the atmosphere}.
\url{https://www.iso.org/standard/17426.html}

\bibitem{iso9613-2}
International Organization for Standardization,
\emph{ISO 9613-2:2024, Acoustics -- Attenuation of sound during propagation
outdoors -- Part 2: Engineering method for the prediction of sound pressure
levels outdoors}.
\url{https://www.iso.org/standard/74047.html}

\bibitem{moller-trumbore}
T. M\"oller and B. Trumbore,
\emph{Fast, Minimum Storage Ray--Triangle Intersection},
Journal of Graphics Tools, 2(1), 21--28, 1997.
\url{https://doi.org/10.1080/10867651.1997.10487468}

\end{thebibliography}

\end{document}
